% =====================================================================
%  FICHE PÉDAGOGIQUE — préambule commun (charte graphique du manuel)
%  Compilation : xelatex (deux passes). Format A4 paysage.
%  Le fichier de contenu définit, avant \begin{document} :
%   \Matiere, \Partie, \NumChap, \TitreChap, \TitreCourt, \Niveau, \Duree, \Banniere, \BanniereDecalage
% =====================================================================
\documentclass[10pt]{article}
\usepackage[a4paper,landscape,top=1.55cm,bottom=1.55cm,left=1.35cm,right=1.35cm,headheight=14pt,headsep=4.5mm,footskip=9mm]{geometry}
\usepackage{fontspec}
\usepackage{polyglossia}
\setmainlanguage{french}
\usepackage{amsmath}
\usepackage{unicode-math}
\setmainfont{TeX Gyre Pagella}
\setsansfont{Poppins}[Path=fonts/, Extension=.ttf, UprightFont=*-Regular, BoldFont=*-Bold, ItalicFont=*-Italic, BoldItalicFont=*-BoldItalic]
\newfontfamily\titlefont{Poppins}[Path=fonts/, Extension=.ttf, UprightFont=*-Bold, BoldFont=*-Bold]
\setmathfont{Latin Modern Math}
\setmathfont[range={it/{latin,Latin,greek,Greek}}]{texgyrepagella-italic.otf}
\setmathfont[range={up/{latin,Latin,num}}]{texgyrepagella-regular.otf}
\usepackage[table]{xcolor}
\usepackage{graphicx,tikz,array,xltabular,multirow}
\usetikzlibrary{calc,fadings}
\usepackage[most]{tcolorbox}
\graphicspath{{images/}}
\usepackage{enumitem}
\usepackage{needspace}
\usepackage{microtype}
\usepackage[newcommands]{ragged2e}
\setlength{\RaggedRightRightskip}{0pt plus 6em}
\usepackage{siunitx}
\sisetup{locale=FR, per-mode=power, inter-unit-product=\ensuremath{\mathord{\cdot}}, group-minimum-digits=5}
\usepackage{fancyhdr}
\setlength{\parindent}{0pt}
\relpenalty=10000 \binoppenalty=10000
\emergencystretch=3em

% ---------- Couleurs de la charte ----------
\definecolor{cInk}{HTML}{1B2A41}  \definecolor{cInk2}{HTML}{51627A}
\definecolor{cBlue}{HTML}{1F5FA8} \definecolor{cOrange}{HTML}{F07F1E}
\definecolor{cRed}{HTML}{D62839}  \definecolor{cGreen}{HTML}{1F9D74}
\definecolor{cPurple}{HTML}{6C4AB6} \definecolor{cTeal}{HTML}{0E8F8A}
\definecolor{cAmber}{HTML}{F4A300} \definecolor{cCyan}{HTML}{0B84C6}
\definecolor{cGrey}{HTML}{8A96A8}  \definecolor{cSoft}{HTML}{F4F7FB}
\color{cInk}

% ---------- Petites commandes ----------
\newcommand\og{«\nobreak\thinspace}
\newcommand\fg{\nobreak\thinspace»}
\newcommand{\sfb}[1]{{\sffamily\bfseries #1}}
\newcommand{\ch}[1]{\ensuremath{\mathrm{#1}}}
\newcommand{\arc}[1]{\overparen{#1}}
% puce de liste dans les cellules : \pt{texte}
\newcommand{\pt}[1]{\par\noindent\hangindent=3.4mm\hangafter=1\makebox[3.4mm][l]{\textcolor{cOrange}{\scriptsize$\blacktriangleright$}}#1}
% sous-puce
\newcommand{\spt}[1]{\par\noindent\hangindent=6.6mm\hangafter=1\hspace*{3.4mm}\makebox[3.2mm][l]{\textcolor{cInk2}{\scriptsize$\bullet$}}#1}
% intertitre dans une cellule
\newcommand{\ct}[2][cInk]{\par\smallskip{\sffamily\bfseries\small\color{#1}#2}\par\nopagebreak\vspace{.5mm}}
% pastille de rubrique : \pastille{couleur}{icône}{texte}
\NewDocumentCommand{\pastille}{m m m}{\tikz[baseline=(t.base)]{\node[fill=#1,text=white,rounded corners=2.6mm,inner xsep=2mm,inner ysep=1.1mm,font=\sffamily\bfseries\scriptsize,anchor=base west](t){\raisebox{-.9mm}{\includegraphics[height=3.6mm]{#2}}\hspace{1.2mm}#3};}}
% étiquette « Diapo n » / « Livre » / « Fiche »
\NewDocumentCommand{\diapo}{m}{\tikz[baseline=(b.base)]\node[fill=cInk,text=white,rounded corners=3pt,inner xsep=2.6pt,inner ysep=1.2pt,font=\sffamily\bfseries\tiny](b){Diapo #1};}
\NewDocumentCommand{\support}{O{cInk2} m}{\tikz[baseline=(b.base)]\node[draw=#1,text=#1,rounded corners=3pt,inner xsep=2.6pt,inner ysep=1.2pt,line width=.5pt,font=\sffamily\bfseries\tiny](b){#2};}
% type d'évaluation
\NewDocumentCommand{\evtype}{m m}{\tikz[baseline=(b.base)]\node[fill=#1!14,text=#1,rounded corners=3pt,inner xsep=3pt,inner ysep=1.3pt,font=\sffamily\bfseries\tiny](b){\MakeUppercase{#2}};\par\vspace{.6mm}}
\newcommand{\diag}{\evtype{cCyan}{Diagnostique}}
\newcommand{\form}{\evtype{cBlue}{Formative}}
\newcommand{\somm}{\evtype{cRed}{Sommative}}
\newcommand{\regul}{\evtype{cGreen}{Remédiation}}
% durée : \duree{min}{cumul}
\newcommand{\duree}[2]{\centering{\titlefont\fontsize{15}{16}\selectfont\color{cInk}#1}\par\vspace{-.4mm}{\sffamily\tiny\color{cInk2}min}\par\vspace{1.4mm}{\sffamily\tiny\color{cGrey}\(\to\) #2}\par}

% ---------- Titres des grandes parties ----------
\newcommand{\bloc}[3]{% \bloc{numéro}{icône}{titre}
  \par\vspace{4mm}\noindent
  \tikz[baseline=(n.base)]\node[fill=cOrange,rounded corners=3pt,text=white,minimum width=9mm,minimum height=9mm,inner sep=0pt,font=\titlefont\fontsize{16}{18}\selectfont](n){#1};%
  \hspace{3mm}\raisebox{-1.8mm}{\includegraphics[height=8mm]{#2}}\hspace{2.5mm}{\titlefont\fontsize{16}{20}\selectfont\color{cInk}#3}\par\vspace{1.3mm}
  {\color{cOrange}\rule{\linewidth}{1.2pt}}\par\vspace{3mm}}

% ---------- Bandeau de la première page ----------
\newcommand{\bandeau}{%
  \thispagestyle{premiere}%
  \begin{tikzpicture}[remember picture,overlay]
    \begin{scope}
      \clip ([yshift=0cm]current page.north west) rectangle ([yshift=-5.1cm]current page.north east);
      \node[anchor=north west,inner sep=0] at ([yshift=\BanniereDecalage]current page.north west) {\includegraphics[width=\paperwidth]{\Banniere}};
      \fill[cInk,path fading=east,opacity=.78] (current page.north west) rectangle ([xshift=19cm,yshift=-5.1cm]current page.north west);
    \end{scope}
    \fill[white] ([yshift=-5.1cm]current page.north west) .. controls ([xshift=10cm,yshift=-4.3cm]current page.north west) and ([xshift=-12cm,yshift=-5.6cm]current page.north east) .. ([yshift=-4.6cm]current page.north east) -- ([yshift=-5.3cm]current page.north east) -- ([yshift=-5.3cm]current page.north west) -- cycle;
    \node[anchor=north west,text width=19.5cm,align=left,inner sep=0] at ([xshift=1.45cm,yshift=-1.0cm]current page.north west) {%
      {\sffamily\bfseries\footnotesize\color{cAmber}\addfontfeatures{LetterSpace=12}FICHE PÉDAGOGIQUE\enspace·\enspace\MakeUppercase{\Matiere}\enspace·\enspace\MakeUppercase{\Partie}}\par\vspace{3.2mm}
      \tikz\node[fill=cOrange,text=white,rounded corners=3.2mm,inner xsep=3.4mm,inner ysep=1.3mm,font=\sffamily\bfseries\normalsize]{Chapitre \NumChap};\par\vspace{2.6mm}
      {\titlefont\fontsize{21}{26}\selectfont\color{white}\TitreChap\par}};
  \end{tikzpicture}%
  \vspace*{3.25cm}\par}

% carte d'information : \infocarte{icône}{couleur}{libellé}{valeur}
\newcommand{\infocarte}[4]{%
  \begin{tcolorbox}[enhanced,width=\linewidth,colback=#2!7!white,colframe=#2!7!white,arc=3mm,boxrule=0pt,left=14mm,right=3mm,top=2mm,bottom=2mm,
    borderline west={2pt}{0pt}{#2},
    overlay={\node at ([xshift=7.2mm]frame.west) {\includegraphics[width=8.4mm]{#1}};}]
  {\sffamily\bfseries\scriptsize\color{#2}\MakeUppercase{#3}}\par\vspace{.4mm}
  {\sffamily\small\color{cInk}\RaggedRight #4\par}
  \end{tcolorbox}}

% ---------- Tableau 1 : libellé de ligne ----------
\newcommand{\lib}[3]{% \lib{icône}{couleur}{libellé}
  \cellcolor{#2!10}\begin{minipage}[t]{\linewidth}\vspace{.6mm}\centering
  \includegraphics[height=10mm]{#1}\par\vspace{1.4mm}
  {\sffamily\bfseries\small\color{#2}#3}\par\vspace{1mm}\end{minipage}}

% ---------- Tableau 2 : bandeau de partie ----------
\newcommand{\partie}[3]{% \partie{couleur}{titre}{durée totale}
  \multicolumn{5}{@{}l@{}}{\cellcolor{#1}\rule{0pt}{6.2mm}\hspace{2.5mm}{\sffamily\bfseries\small\color{white}#2}\hfill{\sffamily\bfseries\scriptsize\color{white}#3}\hspace{2.5mm}}\\[-.2mm]}

% ---------- En-têtes et pieds de page ----------
\pagestyle{fancy}
\fancyhf{}
\fancyhead[L]{\sffamily\footnotesize\color{cInk2}\textcolor{cOrange}{\textbf{FICHE PÉDAGOGIQUE}}\enspace·\enspace\Matiere\enspace·\enspace Chapitre \NumChap}
\fancyhead[R]{\sffamily\footnotesize\color{cInk2}\TitreCourt}
\renewcommand{\headrulewidth}{0.6pt}
\renewcommand{\headrule}{{\color{cBlue!25}\hrule height\headrulewidth width\headwidth}}
\fancyfoot[C]{\tikz\node[circle,fill=cBlue,text=white,font=\sffamily\bfseries\small,minimum size=7mm,inner sep=0pt]{\thepage};}
\fancyfoot[R]{\sffamily\scriptsize\color{cGrey}Durée : \Duree\enspace·\enspace \Niveau}
\fancypagestyle{premiere}{\fancyhf{}\renewcommand{\headrulewidth}{0pt}\fancyfoot[C]{\tikz\node[circle,fill=cBlue,text=white,font=\sffamily\bfseries\small,minimum size=7mm,inner sep=0pt]{\thepage};}}

% ---------- Tableaux ----------
\newcolumntype{L}[1]{>{\RaggedRight\arraybackslash}p{#1}}
\newcolumntype{Y}{>{\RaggedRight\arraybackslash}X}
\newcolumntype{C}[1]{>{\centering\arraybackslash}p{#1}}
\setlength{\arrayrulewidth}{.6pt}
\arrayrulecolor{white}
\keepXColumns
